\section{Discussion}
\label{sec:discussion}

\paragraph{What the argument is, and is not} The claim is not that this card, this model,
or this trace is representative. It is that once sessions can be suspended, admission stops
being a forecast and becomes a feasibility test with a price, and that this reframing is
what turns an unmanaged pool into a usable server under working-set overcommitment. The
quantities in which the argument is stated---bytes of KV, bytes per second between tiers,
rounds per session, rounds per second of prefill---are measurable on any deployment; the
constants are not portable, and we do not present them as such.

\begin{table}[t]
\caption{The regime the design targets, and where it does not apply.}
\label{tab:scope}
\centering
\small
\begin{tabular}{@{}p{0.44\columnwidth}p{0.46\columnwidth}@{}}
\toprule
Applies & Does not apply \\
\midrule
working set $\gg$ KV pool: long sessions, large contexts &
working set well below the pool: nothing is suspended and the ledger is inert \\
single node with a host tier that can hold suspended state &
cluster-level disaggregation, where placement is a different problem \\
contexts large enough that recomputation dominates a copy &
contexts short enough that the recompute is cheaper than the revival \\
sessions that return: multi-round jobs with reusable prefixes &
one-shot requests, where there is no state to preserve \\
\bottomrule
\end{tabular}
\end{table}

\paragraph{Why the negative cases matter} A design of this kind is easy to oversell,
because every mechanism in it can be made to look good on the workload it was tuned for.
The three losing regimes in \S\ref{sec:eval-neg} are therefore part of the result: they
delineate the ratio of working set to capacity at which accounting stops being the
interesting question, and they explain why threshold controllers are not simply mistakes---
in a regime where the working set is small, a window is a perfectly good answer, and the
ledger has nothing to add.

\paragraph{Implications for systems that already have a hierarchy} The mechanism does not
require a new engine. It requires that the engine expose enough for a controller to keep
books---token usage, running counts, per-request cache accounting---and that the scheduler
allow a request to be held at a session boundary rather than only admitted or dropped.
Both are interface changes, not algorithmic ones, and both are absent from current engines
mainly because the request, not the session, is their unit of work.

\paragraph{Threats beyond the evaluation} Our replay arrives all at once, which makes
admission hardest and makes a batch controller look good; a Poisson arrival process would
give the request-level arm more slack and shrink the measured gap, though it would not
change which decisions are feasible. The trace's tool waits are simulated from the
recorded distribution rather than observed live, so the interaction between admission
delays and agent behaviour (an agent that waits longer may emit a longer next round) is
not modelled. Both are directions for measurement rather than for mechanism.
